\section{Fazit}\label{sec:fazit}
Die Schnorr-Identifikation beantwortet die Ausgangsfrage durch das Zusammenspiel von Commitment, unvorhersehbarer Challenge und schlüsselabhängiger Response. Sie ermöglicht einen effizienten Wissensnachweis, ohne den geheimen Schlüssel zu übertragen. Das Zahlenbeispiel zeigt die Funktionsweise vollständig; die Aufwandsabschätzung verdeutlicht, warum große und strukturell geeignete Gruppen nötig sind. Praktische Sicherheit entsteht jedoch erst zusammen mit korrekter Implementierung, frischen Nonces und sicherer Einbettung in ein Authentifizierungssystem. Schnorr ist damit ein vielseitiger kryptografischer Baustein, aber kein vollständiges Identitätsmanagementsystem. Die dargestellten Rechnungen ersetzen weder Messungen konkreter Implementierungen noch eine umfassende Analyse aller aktiven Angriffe.
